\section{Introduction and scope}
\label{sec:scope}

The practical question behind this repository is simple to state. Can one
laptop, with no human in the turn loop, play a current VGC format on the
public Showdown ladder, improve from its own experience, and prove each
improvement before relying on it? Answering it takes several kinds of
component, each with its own evidence and failure modes. A transport has to
deliver legal choices on time. A decision procedure has to choose them. A data
recorder has to preserve exactly what was known when each choice was made. A
learner has to turn those records into candidate models. A statistical gate
has to decide which candidates are better. This report explains each
component, gives the mathematics it relies on, and states what has been
measured.

\subsection{What changed since the first edition}
The first edition (7 October 2026) described a persistent MCP harness in
which the OpenClaw agent called tools and a local actor--critic chose every
battle decision. Thirty-three public commits later the deployed system is
different in five ways that matter mathematically:
\begin{enumerate}
  \item \textbf{Turn decisions come from search.} A determinized one-turn
    simultaneous-move search over the official simulator chooses live turn
    actions (Section~\ref{sec:search}). The neural actor now chooses only team
    previews, and it does so greedily.
  \item \textbf{The neural policy has a temperature and a strategic prior.}
    Logits are divided by $\tau=0.25$. The additive prior is no longer the
    legacy power heuristic: it is a squashed one-turn scenario evaluation
    (Sections~\ref{sec:features} and~\ref{sec:network}).
  \item \textbf{Promotion is a hypothesis test.} Candidates must pass an exact
    one-sided paired sign test at $p\le0.05$ as well as a win margin. The worker
    only stages a candidate; promotion happens later at a matchmaking boundary
    (Section~\ref{sec:promotion}).
  \item \textbf{Training is a three-lane pipeline.} CPU collection, PPO on the
    Apple~MPS GPU and CPU evaluation run as separate bounded processes. Each
    trajectory is checked against its collecting likelihood, and PPO stops
    early on a KL criterion (Sections~\ref{sec:learning} and~\ref{sec:pipeline}).
  \item \textbf{Live play is browser-driven.} Playwright clicks the official
    client, guarded by the request identifier. OpenClaw and MCP remain available
    for agent-supervised play but are no longer in the turn loop
    (Section~\ref{sec:transport}).
\end{enumerate}
Appendix~\ref{app:errata} lists every first-edition statement that this
edition corrects, including one circular proposition and several claims made
stale by later commits.

\subsection{Three kinds of statement}
The report keeps three kinds of statement apart:
\begin{enumerate}
  \item \textbf{Implemented behaviour}: facts traceable to source functions,
    such as tensor dimensions, constants and gate predicates. These are
    followed by an \emph{Implementation} line naming the files.
  \item \textbf{Conditional mathematical results}: propositions proved under
    explicit assumptions, followed by an \emph{Interpretation} or
    \emph{Caveat} that says how far the result transfers to the running system.
  \item \textbf{Measurements}: results of experiments run by the project,
    cited to the repository's validation documents or, where marked, to the
    local run ledger. This paper ran no new games, training or promotion, apart
    from read-only recomputation of the statistics it reports.
\end{enumerate}

\subsection{Format and pinned mechanics}
The main target is \code{gen9championsvgc2026regmc}: doubles, bring six and
pick four, open team sheets on request, one Mega Evolution per battle, no
Terastallization, Item and Species Clauses, and Champions stat points instead
of EVs and IVs \citep{championsrules}. The local simulator dependency is
pinned to official Showdown commit
\code{c046106cbe075931b1ff8d8b800ff5be47a85f96} \citep{showdownsource}. A
format identifier beginning with \code{gen9} does not determine these
mechanics. Format, simulator revision, feature profile, team fingerprint and
collecting checkpoint are part of every scientific statement in this report.

\subsection{Summary of measured results}
\begin{table}[htbp]
\centering\small
\begin{tabularx}{\linewidth}{>{\raggedright\arraybackslash}p{0.34\linewidth} X}
\toprule
Question & Measured answer (details in Section~\ref{sec:evidence}) \\
\midrule
Does search beat the learned policy offline? & Same 45 seeds, sides and real
opponent teams: search 30/45 vs.\ policy 19/45; 14 gains, 3 losses; exact
McNemar $p=0.0064$. Clustering by opponent team: $p=0.0037$. \\
Did the learned policy improve over campaigns? & Live record 348--499 (41.1\%)
over 847 games across three actor generations. One strategic prior passed an
independent 160-game final (125 vs.\ 106, $p=0.0072$). One PPO update passed
the 60-game gate out of 129 attempts. \\
What happened live after search? & Preliminary and observational: 41 games,
27--14. Rating 1182 $\rightarrow$ 1464. Elo-expected 18.4 wins in 37 rated
games; observed 23 (exact $p=0.086$). \\
What failed? & More than thirty candidates: PPO variants, preview models, a
matchup architecture that passed development and then tied its final, and an
offline hyperparameter matrix with no significant arm after Holm correction. \\
\bottomrule
\end{tabularx}
\caption{Headline evidence. Ladder results depend on the opponent pool and on
rating-based matchmaking; Section~\ref{sec:elo} explains why a raw win rate is
not a stable strength measure.}
\label{tab:headline}
\end{table}

\subsection{How to read this report}
Section~\ref{sec:system} gives the system map and project timeline.
Sections~\ref{sec:game}--\ref{sec:network} formalize the game, the candidate
actions, the features and the neural policy. Section~\ref{sec:data} describes
the record kept for every match. Sections~\ref{sec:learning}--\ref{sec:promotion}
cover learning, the CPU/MPS pipeline and promotion statistics.
Section~\ref{sec:search} analyses the search agent and
Section~\ref{sec:transport} the live turn. Section~\ref{sec:scout} covers
opponent-move scouting and team selection. Section~\ref{sec:evidence} collects
all evidence and Section~\ref{sec:limits} its limits. The appendices provide a
glossary, the corrections to the first edition, defects found during review, a
source map and reproduction notes.
